\documentclass[12pt,a4paper]{article}
\usepackage[russian]{babel}
\usepackage{fontspec}
\usepackage{geometry}
\usepackage{booktabs}
\usepackage{longtable}
\usepackage{array}
\usepackage{multirow}
\usepackage{enumitem}
\usepackage{titlesec}
\usepackage{xcolor}
\usepackage{siunitx}

\usepackage{pdfpages}

\usepackage{tablefootnote}

% Дополнительные пакеты для данного документа
\usepackage{caption}
\usepackage{subcaption}
\usepackage{listings}
% \usepackage{hyperref}

\geometry{a4paper, margin=25mm}
\setmainfont{Liberation Serif}
\setsansfont{Liberation Sans}

\definecolor{adblue}{RGB}{0,84,166}

\titleformat{\section}{\large\bfseries\color{adblue}}{}{0em}{}[\titlerule]
\titleformat{\subsection}{\normalsize\bfseries}{}{0em}{}
\titleformat{\subsubsection}{\normalsize\itshape}{}{0em}{}

\sisetup{output-decimal-marker={,}}

\usepackage{tikz}
\usetikzlibrary{positioning, calc, backgrounds, math, 3d, perspective, arrows.meta, graphs, graphdrawing, fit, shapes.geometric, trees}
    

\usepackage{pgfplots} % для построения графиков
\pgfplotsset{compat=1.18} % версия для совместимости
\usepackage{tkz-euclide}

% --- Подключение пакета алгоритмов (без флага algo2e, чтобы вернуть \begin{algorithm}) ---
\usepackage[ruled,vlined]{algorithm2e}
%\usepackage{caption}
\newcommand{\Gray}[1]{\textcolor{gray}{#1}} 
% Настройка подписей
\DeclareCaptionLabelFormat{gost}{#1 #2}
\captionsetup[table]{labelformat=gost,labelsep=endash,justification=justified,singlelinecheck=false,font=normal}
\captionsetup[figure]{labelformat=gost,labelsep=endash,justification=centering,font=normal}

% Локализация algorithm2e на русский (ОБЯЗАТЕЛЬНО после загрузки пакета)
\SetAlgorithmName{Алгоритм}{}{}
\SetAlgoCaptionSeparator{---}
\SetKwInput{Input}{Вход}
\SetKwInput{Output}{Выход}

\usepackage{url}
% Говорим пакету использовать моноширинный шрифт (как у texttt)
\Urlmuskip=0mu plus 1mu % гибкие пробелы для формул, если нужно
\DeclareUrlCommand\code{\urlstyle{tt}} 

% Библиотека схемотехники
\usepackage[siunitx, RPvoltages, american]{circuitikzgit}
\ctikzset{
    amplifiers/fill=cyan,
    sources/fill=green!20,
    diodes/fill=white,
    resistors/fill=violet,
    grounds/scale=1.2,
}

\usepackage{iftex}

\ifXeTeX % Если используется TeTeX или TeLaTeX
  \typeout{Режим XeTeX}
  \usepackage{xltxtra}
  \usepackage{fontspec}
  \usepackage{polyglossia}
  \setmainlanguage{russian}
  \setotherlanguage{english}
  \setkeys{russian}{babelshorthands=true} % По идее включает <<, >>, -- ---

  \setmainfont{Times New Roman}[Ligatures=TeX]
  \setromanfont{Times New Roman}[Ligatures=TeX]
  \setsansfont{Arial}[Ligatures=TeX]
  \setmonofont{Courier New}[Ligatures=TeX] 

  \newfontfamily{\cyrillicfont}{Times New Roman}
  \newfontfamily{\cyrillicfontrm}{Times New Roman}
  \newfontfamily{\cyrillicfonttt}{Courier New}
  \newfontfamily{\cyrillicfontsf}{Arial}
\fi

\ifLuaTeX % Если используется LuaLatex
  \typeout{Режим LuaTeX}
%  \usepackage{luatextra}
  \usepackage{fontspec}
  \defaultfontfeatures{Ligatures={TeX}}
  \usepackage{polyglossia}
  \setmainlanguage{russian}
  \setotherlanguage{english}
  \setmainfont{Times New Roman}[Script=Cyrillic, Ligatures=TeX]
  \setromanfont{Times New Roman}[Script=Cyrillic, Ligatures=TeX]
  \setsansfont{Arial}[Script=Cyrillic, Ligatures=TeX]
  \setmonofont{Courier New}[Ligatures=TeX]
  \newfontfamily{\cyrillicfont}{Times New Roman}[Script=Cyrillic, Ligatures=TeX]
  \newfontfamily{\cyrillicfontrm}{Times New Roman}[Script=Cyrillic, Ligatures=TeX]
  \newfontfamily{\cyrillicfonttt}{Courier New}
  \newfontfamily{\cyrillicfontsf}{Arial}[Script=Cyrillic, Ligatures=TeX]

  \usegdlibrary{trees}
  \usegdlibrary{force}
\fi

\ifPDFTeX % Если используется 8битный PDFTex
  \typeout{8-битный режим pdfTex}
  % For pdfTeX we must use old
  % 8-bit font technologies
  \usepackage[T2A]{fontenc}
  \usepackage[english,russian]{babel}
  \usepackage[utf8]{inputenc}
  \usepackage[english, russian]{babel}

  %\usepackage{newtxmath}
  \usepackage{amsmath}
  \usepackage{amsthm}
  \usepackage{amssymb}
  \usepackage{amsfonts}
  \usepackage{mathtext}
\fi

\usepackage{amsthm}  % Возможности AMSMath
\usepackage{amsmath}
\usepackage{amssymb}

\usepackage{esvect} % Для красивых стрелок
\usepackage{bm} % Для жирных символов

%Hyphenation rules
\usepackage{hyphenat}
\hyphenation{ма-те-ма-ти-ка вос-ста-нав-ли-вать ото-бра-жа-ют-ся сге-не-ри-ро-ван-ном сим-во-лы}
% \usepackage[english, russian]{babel}

\usepackage{graphicx}
\graphicspath{ {./images/} }

\usepackage{empheq}

\usepackage{float}

% \usepackage{comment}

\usepackage[hidelinks,
    unicode=true,          % поддерживает Unicode в закладках
    psdextra=true,      % автоматически конвертирует простую математику в текст
    % focusdefault=false
]{hyperref}


% Окружение "Теорема"
\newtheorem{theorem}{Теорема}[section]
\newtheorem{corollary}{Corollary}[theorem]
\newtheorem{lemma}[theorem]{Lemma}

% Окружение "Определение"
\theoremstyle{definition} % Прямой шрифт, нумерация
\newtheorem{definition}{Определение}[section]
\newtheorem{example}{Пример}[section]

% Окружение "Замечание"
\theoremstyle{remark} % Прямой шрифт, без нумерации
% Окуржение "Решение"
\newtheorem*{solution}{Решение}
\newtheorem*{remark}{Замечание}

\usepackage{polynom}
\usepackage{tabularx}
\usepackage{cancel}

\addto\captionsrussian{%
  \renewcommand{\figurename}{Рис.}%
  \renewcommand{\tablename}{Табл.}%
}

\renewcommand{\arraystretch}{1.5} % Высота строки таблицы

\renewcommand{\labelitemi}{\(\circ\)} % Бюллет -- пустой кружок в itemize

\title{Список характеристик, рассчитываемых во-время обучения модели}
\author{Наталия А. Рыбкина}
\date{2026}

\begin{document}

\maketitle

\section{Функции потерь и их роль в обучении}

\subsection{Пространственная функция потерь L1Loss}

\subsubsection{Простое объяснение}

В NAFNet (и других сетях реставрации изображений) — это обычное среднее арифметическое абсолютных отклонений по каждому пикселю, где у каждого пикселя абсолютно одинаковый (равный) статический вес. 

В стандартной реализации L1Loss (как и MSELoss) каждый пиксель вносит вклад с одинаковым статическим весом, равным 1/N, где N – общее количество пикселей. Формула, которую вы привели:

\[
\mathcal{L}_1 = \frac{1}{N}\sum_{i=1}^N |I_{\text{pred}}(i) - I_{\text{gt}}(i)|
\]

означает, что \textbf{все пространственные позиции (и все цветовые каналы) имеют равную важность} для оптимизации. Это основное ограничение таких «пиксельных» функций потерь: они не учитывают структуру изображения, восприятие человеком или частотные характеристики.  

Именно поэтому предлагается добавить \textbf{комбинированную потерю} с частотной компонентой (FocalFrequencyLoss), которая:

\begin{itemize}
    \item динамически изменяет вес различных частотных компонент (а не пикселей) в зависимости от того, насколько сильно сеть ошибается на каждой частоте;
    \item тем самым заставляет модель восстанавливать высокочастотные детали (границы, текстуры), которые при использовании одного L1Loss часто сглаживаются.
\end{itemize}

В коде nn.L1Loss() вычисляется среднее арифметическое абсолютных отклонений по всем элементам тензора (по измерениям: батч, каналы, высота, ширина). Веса не меняются в процессе обучения, не зависят от пространственного положения или цвета пикселя.

\subsubsection{Подробно}

В задачах восстановления изображений одной из базовых функций потерь является \textbf{L1Loss} (средняя абсолютная ошибка). Для пары изображений — предсказанного сетью \(I_{\text{pred}}\) и эталонного \(I_{\text{gt}}\) — она определяется как
\[
\mathcal{L}_1 = \frac{1}{N}\sum_{i=1}^{N} \bigl|I_{\text{pred}}(i) - I_{\text{gt}}(i)\bigr|,
\]
где \(N\) — общее число пикселей. L1Loss обеспечивает устойчивую сходимость и базовую геометрическую точность, однако обладает существенными недостатками применительно к задаче демозаики и шумоподавления: она не различает низко- и высокочастотные искажения, способствует сглаживанию резких границ («эффект мыла») и не подавляет цветовой муар, возникающий из-за наложения спектров при дискретизации.

\subsection{Комбинированная частотно-пространственная потеря (CombinedLoss)}

Для преодоления ограничений чисто пространственных метрик в данной работе предложена комбинированная функция потерь
\[
\mathcal{L}_{\text{total}} = \gamma_1 \cdot \mathcal{L}_1 + \gamma_2 \cdot \mathcal{L}_{\text{FFL}},
\]
где \(\mathcal{L}_{\text{FFL}}\) — \textbf{Focal Frequency Loss} (частотная потеря с фокусировкой). Её вычисление включает следующие шаги:
\begin{enumerate}
    \item Двумерное быстрое преобразование Фурье (БПФ) предсказания \(I_{\text{pred}}\) и эталона \(I_{\text{gt}}\).
    \item Вычисление амплитудных спектров \(A_{\text{pred}}(u,v)\) и \(A_{\text{gt}}(u,v)\).
    \item Определение квадрата разности амплитуд \(\Delta A(u,v) = \bigl(A_{\text{pred}}(u,v) - A_{\text{gt}}(u,v)\bigr)^2\).
    \item Построение динамических весов \(w(u,v) = \bigl(\Delta A(u,v) / \max(\Delta A)\bigr)^{\alpha}\), где \(\alpha\) — параметр фокусировки.
    \item Итоговая частотная потеря:
    \[
    \mathcal{L}_{\text{FFL}} = \frac{1}{HW}\sum_{u,v} w(u,v) \cdot \Delta A(u,v).
    \]
\end{enumerate}
Комбинированная потеря позволяет одновременно минимизировать пространственные ошибки и восстанавливать высокочастотные компоненты, подавляя артефакты алиасинга и муара.

\subsection{Периодичность вычисления и логирования}

\paragraph{Обучение (тренировочные метрики).}
Loss (как \(\mathcal{L}_1\), так и \(\mathcal{L}_{\text{total}}\)) вычисляется на \textbf{каждой итерации} (каждом батче) в процессе прямого распространения сигнала. Значение loss используется для обратного распространения ошибки и обновления весов модели. В дополнение к loss на каждой итерации также вычисляются вспомогательные метрики:
\begin{itemize}
    \item PSNR (пиковое отношение сигнал/шум) на тренировочном батче,
    \item текущая скорость обучения (learning rate),
    \item изменение PSNR относительно предыдущего шага (\(\Delta\)PSNR),
    \item дисперсия градиентов,
    \item отношение полных вариаций (TV ratio),
    \item использование видеопамяти.
\end{itemize}
Все эти метрики записываются в CSV-файл (по умолчанию \texttt{train\_metrics.csv}) на каждом шаге, что обеспечивает возможность детального анализа сходимости.

\paragraph{Валидация (тестовые метрики).}
Метрики, требующие прогона всей тестовой выборки, такие как структурное сходство (SSIM) и, при необходимости, валидационный PSNR, вычисляются с заданной в конфигурационном файле периодичностью — каждые \(k\) эпох (параметр \texttt{validation\_freq}). Эти значения сохраняются либо в отдельный CSV-файл (\texttt{val\_metrics.csv}), либо могут быть извлечены из общего лога. Такая редкая периодичность обусловлена высокими вычислительными затратами на обработку полной тестовой выборки.

\subsection{Практическая реализация}

В разработанном программном комплексе класс \texttt{CombinedLoss} уже реализован в модуле \texttt{libraries/training\_losses.py}. Для его активации в процессе обучения достаточно в скрипте \texttt{run\_training.py} заменить строку
\begin{verbatim}
criterion = nn.L1Loss()
\end{verbatim}
на
\begin{verbatim}
criterion = CombinedLoss(opt)
\end{verbatim}
а в YAML-конфигурации добавить секцию \texttt{losses} с весами \(\gamma_1\) и \(\gamma_2\). Рекомендуемые значения: \(\gamma_1 = 1.0\), \(\gamma_2 = 0.5\ldots 1.0\), \(\alpha = 1.0\). Параметр \texttt{validation\_freq} задаётся в секции \texttt{train} конфига и определяет, как часто выполняются валидационные расчёты (например, \texttt{validation\_freq: 1} означает каждую эпоху).

\subsection{Комбинированная потеря (CombinedLoss)}

Как было показано ранее, использование только пиксельной функции потерь (L1Loss) приводит к сглаживанию высокочастотных деталей и не позволяет эффективно бороться с цветовым муаром, возникающим при демозаике. Для преодоления этих ограничений в диссертации предложена и реализована комбинированная функция потерь, объединяющая пространственный и частотный критерии:

\[
\mathcal{L}_{\text{total}} = \gamma_1 \cdot \mathcal{L}_1 + \gamma_2 \cdot \mathcal{L}_{\text{FFL}}.
\]

Здесь \(\mathcal{L}_1\) — стандартная средняя абсолютная ошибка (L1Loss), вычисляемая как
\[
\mathcal{L}_1 = \frac{1}{N}\sum_{i=1}^{N}\bigl|I_{\text{pred}}(i) - I_{\text{gt}}(i)\bigr|,
\]
где \(N\) — общее число пикселей во всех каналах изображения. Данная составляющая обеспечивает грубую геометрическую согласованность и устойчивость обучения.

Второе слагаемое \(\mathcal{L}_{\text{FFL}}\) — \textbf{Focal Frequency Loss} (частотная потеря с фокусировкой). Её вычисление включает следующие этапы:
\begin{enumerate}
    \item Вычисление двумерного дискретного преобразования Фурье (БПФ) для предсказания и эталона:
    \[
    F_{\text{pred}}(u,v) = \mathcal{F}\{I_{\text{pred}}\}(u,v), \qquad 
    F_{\text{gt}}(u,v) = \mathcal{F}\{I_{\text{gt}}\}(u,v).
    \]
    \item Определение амплитудных спектров:
    \[
    A_{\text{pred}}(u,v) = |F_{\text{pred}}(u,v)|, \qquad
    A_{\text{gt}}(u,v) = |F_{\text{gt}}(u,v)|.
    \]
    \item Вычисление квадрата разности амплитуд:
    \[
    \Delta A(u,v) = \bigl(A_{\text{pred}}(u,v) - A_{\text{gt}}(u,v)\bigr)^2.
    \]
    \item Построение динамической весовой матрицы (фокусировка):
    \[
    w(u,v) = \left(\frac{\Delta A(u,v)}{\max(\Delta A) + \varepsilon}\right)^{\alpha},
    \]
    где \(\alpha\) — параметр фокусировки (в работе принят \(\alpha = 1.0\)), \(\varepsilon\) — малая константа для устойчивости.
    \item Итоговая частотная потеря:
    \[
    \mathcal{L}_{\text{FFL}} = \frac{1}{H W}\sum_{u=0}^{H-1}\sum_{v=0}^{W-1} w(u,v) \cdot \Delta A(u,v).
    \]
\end{enumerate}
Благодаря динамическим весам \(w(u,v)\) частотная компонента фокусируется на тех пространственных частотах, где модель допускает наибольшие ошибки в спектре. Это заставляет нейросеть восстанавливать высокочастотные детали (резкие границы, мелкие текстуры), которые при использовании одного \(\mathcal{L}_1\) неизбежно сглаживаются.

Параметры \(\gamma_1\) и \(\gamma_2\) задают баланс между пространственной точностью и качеством частотного восстановления. В экспериментах настоящей работы установлены значения \(\gamma_1 = 1.0\), \(\gamma_2 = 1.0\). При необходимости веса могут быть изменены в конфигурационном файле в секции \texttt{losses}. 

Комбинированная потеря вычисляется на каждом шаге обучения (каждом батче) и используется для обратного распространения ошибки. Её значение логируется в текстовый файл и CSV на каждой итерации (с заданной в конфиге частотой), что позволяет отслеживать сходимость как пространственной, так и частотной составляющих.

\subsection{Пиковое отношение сигнал/шум (PSNR)}

\subsubsection{Простое объяснение}

Почему же \textbf{PSNR} (пиковое отношение сигнал/шум), будучи прекрасной метрикой для оценки качества, почти никогда не используется в качестве функции потерь (целевой функции) при обучении нейросетей восстановления изображений? Разберём «на пальцах» два ключевых отличия PSNR от классической L1Loss.

\paragraph{Отличие 1: квадрат \textit{vs.} модуль.}
\begin{itemize}
    \item \textbf{L1Loss} считает модуль ошибки: \(|I_{\text{pred}} - I_{\text{gt}}|\). Для неё любые ошибки «равны»: один пиксель, ошибочный на 10 единиц яркости, даст такой же штраф, как десять пикселей с ошибкой в 1 единицу (штраф = 10).
    \item \textbf{PSNR} в своей основе использует среднеквадратичную ошибку (MSE): \((I_{\text{pred}} - I_{\text{gt}})^2\). Здесь крупные промахи наказываются непропорционально сильно: ошибка на 10 единиц даёт вклад \(10^2 = 100\), а десять маленьких ошибок — всего \(10\cdot1^2 = 10\).
\end{itemize}
\textbf{Следствие для сети:} если пытаться минимизировать PSNR напрямую, сеть будет \textbf{панически бояться редких, но грубых артефактов} (например, размытия края) и бросит все ресурсы на их устранение, часто в ущерб мелким текстурам. L1 же относится к ошибкам более «справедливо», позволяя сохранять естественную зернистость.

\paragraph{Отличие 2: логарифм и поведение градиента.}
В PSNR присутствует логарифм: \(\text{PSNR}=10\log_{10}(1/\text{MSE})\). Логарифм кардинально меняет динамику градиента на разных этапах обучения:
\begin{itemize}
    \item \textbf{В начале обучения}, когда MSE велико (сеть выдаёт «кашу»), производная логарифма мала. Градиенты почти затухают, сеть обучается \textit{очень медленно}.
    \item \textbf{В конце обучения}, когда MSE становится крошечным (сеть почти идеальна), логарифм резко «взвинчивает» градиенты. Сеть начинает штрафовать даже микроскопические ошибки, которые при L1 уже были бы проигнорированы.
\end{itemize}
\textbf{Следствие для сети:} обучение с PSNR-лоссом нестабильно — вначале оно почти не идёт, а на финише становится чрезмерно чувствительным к шуму. К тому же при \(\text{MSE}=0\) функция терпит разрыв.

\paragraph{Итог.}
В современных архитектурах реставрации \textbf{PSNR используется только как оценочная метрика} (для сравнения с классическими методами и для мониторинга), но не как функция потерь. Вместо неё применяют L1Loss (или её гладкий вариант CharbonnierLoss) для устойчивой сходимости, а для улучшения резкости и восприятия — комбинированные потери (например, CombinedLoss = L1 + Focal Frequency Loss), о которых рассказано выше. Таким образом, отсутствие PSNR в роли лосса — это осознанное инженерное решение, а не упущение.

\subsubsection{Подробно}

\textbf{Пиковое отношение сигнал/шум} (Peak Signal-to-Noise Ratio, \textbf{PSNR}) является одной из наиболее распространённых объективных метрик качества восстановления изображений. Она измеряет отношение максимально возможной мощности сигнала к мощности шума (ошибки) между восстановленным изображением \(I_{\text{pred}}\) и эталонным \(I_{\text{gt}}\).

Для изображений, представленных в диапазоне \([0,1]\) (нормализованные значения), PSNR определяется через среднеквадратичную ошибку (MSE):

\[
\text{MSE} = \frac{1}{N}\sum_{i=1}^{N}\bigl(I_{\text{pred}}(i) - I_{\text{gt}}(i)\bigr)^2,
\]

\[
\text{PSNR} = 10 \cdot \log_{10}\left(\frac{1}{\text{MSE}}\right) \quad \text{(дБ)}.
\]

Здесь \(N\) — общее число пикселей во всех каналах, а максимальная мощность сигнала принята за \(1\). Чем выше PSNR, тем меньше искажений. Значения PSNR порядка \textbf{30–40 дБ} обычно соответствуют хорошему качеству, а \textbf{>40 дБ} — очень высокому.

\paragraph{Достоинства PSNR:}
\begin{itemize}
    \item \textbf{Простота вычисления} и интерпретации.
    \item \textbf{Математическая строгость} — прямо связана с MSE.
    \item \textbf{Широкая распространённость} в литературе по обработке изображений.
\end{itemize}

\paragraph{Ограничения PSNR:}
\begin{itemize}
    \item \textbf{Плохо коррелирует с субъективным зрительным восприятием}: небольшие геометрические искажения или сдвиг цвета могут давать высокий PSNR, но визуально картинка плоха.
    \item \textbf{Не учитывает структурную информацию} (потеря резких границ может слабо отражаться на PSNR).
    \item \textbf{Чувствителен к глобальному сдвигу яркости} и не различает типы артефактов (муар, блочность, размытие).
\end{itemize}

В данной работе PSNR вычисляется на \textbf{каждой итерации} (каждом батче) для тренировочных данных и используется как вспомогательная метрика для мониторинга сходимости. Значения PSNR логируются в CSV-файл на каждом шаге, что позволяет строить кривые обучения. Кроме того, при валидации (с заданной периодичностью) может вычисляться PSNR на тестовой выборке, что даёт более объективную оценку обобщающей способности модели.

Несмотря на недостатки, PSNR остаётся стандартным инструментом для сравнения с классическими методами (билинейная интерполяция, MHC) и является обязательным в техническом задании на диссертацию.

\subsection{Скорость обучения (Learning Rate, LR)}

\subsubsection{Простое объяснение}

\textbf{Скорость обучения} — это размер шага, который делает оптимизатор при движении в сторону минимума функции потерь. Представьте, что вы спускаетесь с горы в темноте: вы делаете шаг, ощупываете склон и решаете, куда идти дальше. Если шаги слишком большие, вы можете перепрыгнуть через яму (минимум) и укатиться в другую сторону. Если шаги крошечные, вы будете спускаться бесконечно долго, рискуя застрять на мелкой кочке. Скорость обучения управляет длиной этих шагов.

В начале обучения, когда сеть ничего не умеет, можно делать относительно большие шаги — это быстро приближает нас к области хороших решений. Постепенно, когда мы приблизились к минимуму, шаги нужно уменьшать, чтобы не «перепрыгнуть» дно и аккуратно «вылизать» оптимальные веса. Именно для этого используют \textbf{планировщики скорости обучения} (schedulers).

\subsubsection{Подробно}

Формула обновления весов для градиентного спуска (в простейшей форме):
\[
w_{t+1} = w_t - \eta \cdot \nabla \mathcal{L}(w_t),
\]
где \(\eta\) — скорость обучения, \(\nabla \mathcal{L}\) — градиент функции потерь. 

\textbf{Влияние слишком большого \(\eta\)}: веса сильно «дрожат», потеря может расходиться или колебаться без уменьшения. \textbf{Слишком малое \(\eta\)}: обучение застревает в локальных минимумах или требует огромного числа итераций.

В данной работе используется алгоритм \textbf{AdamW} (Adam с коррекцией весов). Он автоматически адаптирует скорость обучения для каждого параметра, используя моменты первого и второго порядка. Однако даже для AdamW полезно дополнительно управлять глобальным темпом обучения с помощью планировщика.

Выбранный планировщик — \textbf{CosineAnnealingLR}. Он изменяет \(\eta\) в соответствии с косинусной функцией:
\[
\eta(t) = \eta_{\min} + \frac{1}{2}(\eta_{\max} - \eta_{\min})\left(1 + \cos\left(\frac{t}{T_{\max}}\pi\right)\right),
\]
где \(t\) — текущий номер эпохи (или шага), \(T_{\max}\) — заданное количество эпох, \(\eta_{\max}\) — начальная скорость обучения, \(\eta_{\min}\) — минимальное значение в конце расписания. Такое плавное, «мягкое» снижение позволяет сети сначала быстро улучшаться, а затем тонко подстраивать веса, избегая резких скачков.

\subsubsection{Использование при обучении}

В конфигурационном файле эксперимента задаются следующие параметры, относящиеся к скорости обучения:

\begin{verbatim}
train:
  optim_g:
    lr: 1e-3                     # начальная скорость обучения
    weight_decay: 0              # L2-регуляризация (отключена)
    betas: [0.9, 0.9]            # параметры AdamW
  scheduler:
    eta_min: 1e-7                # минимальная скорость обучения в конце
\end{verbatim}

\begin{itemize}
    \item Начальное значение \(\eta = 10^{-3}\) — типичный выбор для AdamW в задачах реставрации изображений. Оно обеспечивает быструю начальную сходимость без нестабильности.
    \item Планировщик настроен на понижение \(\eta\) до \(10^{-7}\) за \(T_{\max} = \text{num\_epochs}\) эпох. Такое сильное уменьшение гарантирует, что в конце обучения сеть делает очень маленькие шаги, «вылизывая» веса.
    \item \textbf{Причина выбора CosineAnnealingLR}: он не имеет резких изломов (в отличие от StepLR или ReduceLROnPlateau), что даёт предсказуемое поведение и хорошо сочетается с AdamW. Многие современные архитектуры (включая NAFNet) используют именно его.
\end{itemize}

В процессе обучения на каждой итерации вычисляется текущее значение LR (согласно расписанию) и логируется в CSV-файл. Это позволяет впоследствии проанализировать, не случилось ли преждевременного затухания градиентов.

\subsection{Дельта PSNR (\(\Delta\)PSNR) и дисперсия градиентов (GradVar)}

\subsubsection{Простое объяснение}

\textbf{\(\Delta\)PSNR} — это \textit{изменение} значения PSNR между последовательными шагами (итерациями) обучения. Если обучение идёт успешно, PSNR должен расти, а \(\Delta\)PSNR — быть положительным. Когда \(\Delta\)PSNR становится близким к нулю или колеблется около нуля, это сигнал, что модель вышла на плато или сошлась. В диссертации эта метрика используется как ранний индикатор остановки сходимости.

\textbf{Дисперсия градиентов (GradVar)} — это статистический разброс величин градиентов весов по всем параметрам модели. Если градиенты становятся \textit{слишком маленькими} (дисперсия близка к нулю), сеть перестаёт обучаться — это явление называется \textit{затуханием градиентов} (vanishing gradients). Если дисперсия \textit{слишком велика} (градиенты сильно отличаются по модулю), обучение становится хаотичным и нестабильным. В хорошо настроенной сети GradVar обычно медленно уменьшается в процессе обучения, оставаясь в разумных пределах.

\subsubsection{Подробно}

\paragraph{\(\Delta\)PSNR.}
Формально:
\[
\Delta\text{PSNR}(t) = \text{PSNR}(t) - \text{PSNR}(t-1),
\]
где \(t\) — номер итерации (глобальный шаг). \(\Delta\)PSNR вычисляется в \texttt{TrainingLogger} и сохраняется в CSV-файл на каждом шаге. Положительное значение означает улучшение качества, отрицательное — ухудшение (возможно, из-за выхода из локального минимума или слишком высокой скорости обучения). На практике \(\Delta\)PSNR редко используется как самостоятельный критерий остановки, но в совокупности с loss и другими метриками помогает понять динамику обучения.

\paragraph{Дисперсия градиентов (GradVar).}
Пусть \(\mathbf{g} = \{g_1, g_2, \dots, g_M\}\) — все компоненты градиентов всех параметров модели на данной итерации. Тогда
\[
\overline{g} = \frac{1}{M}\sum_{i=1}^{M} g_i, \qquad 
\text{GradVar} = \frac{1}{M}\sum_{i=1}^{M} (g_i - \overline{g})^2.
\]
В коде вычисляется как \texttt{np.var(grad\_norms)}, где \texttt{grad\_norms} — список норм градиентов для каждого слоя (а не для каждого параметра). Это более грубая оценка, но достаточная для мониторинга.

\textbf{Интерпретация:}
\begin{itemize}
    \item \textbf{Слишком малая GradVar (например, \(<10^{-8}\))}: градиенты почти занулились — возможно, сеть попала в плато или перестала обучаться из-за слишком маленькой скорости обучения или проблем с архитектурой.
    \item \textbf{Слишком большая GradVar (например, \(>10^{-2}\))}: градиенты сильно различаются, некоторые параметры обновляются гораздо быстрее других — это может приводить к нестабильности и расходимости. Часто наблюдается в начале обучения при неподходящей скорости обучения.
    \item \textbf{Умеренное, плавно уменьшающееся значение}: признак здорового обучения.
\end{itemize}

\subsubsection{Использование при обучении}

Обе метрики вычисляются в \texttt{TrainingLogger.log\_metrics} на \textbf{каждой итерации} (каждом шаге) и записываются в CSV-файл (\texttt{train\_metrics.csv}) при условии, что соответствующие флаги в YAML-конфиге включены (по умолчанию \texttt{true} для \texttt{include\_in\_file\_log}). Вывод в консоль регулируется отдельно.

В диссертации эти метрики не являются центральными, но помогают:
\begin{itemize}
    \item \textbf{Для \(\Delta\)PSNR}: убедиться, что PSNR не просто стабилизировался, а действительно достиг предела (когда \(\Delta\)PSNR становится шумом около нуля).
    \item \textbf{Для GradVar}: диагностировать возможное затухание градиентов (особенно важно для глубоких сетей, таких как NAFNet с 12 средними блоками). Если GradVar резко падает, можно скорректировать скорость обучения или архитектуру.
\end{itemize}

На практике в данной работе \(\Delta\)PSNR и GradVar используются только для внутреннего мониторинга и не влияют на процесс обучения (не участвуют в принятии решений об остановке или изменении гиперпараметров). Тем не менее, их наличие в логах облегчает отладку и тонкую настройку эксперимента.

\subsection{Дисперсия градиентов (GradVar)}

\subsubsection{Простое объяснение}

\textbf{Дисперсия градиентов} — это мера разброса величин градиентов по всем параметрам нейросети. Представьте, что вы учите группу студентов решать одну и ту же задачу. Если все делают ошибки примерно одинакового размера, дисперсия мала. Если кто-то ошибается сильно, а кто-то почти идеален, дисперсия велика.

В обучении нейросетей градиенты указывают, насколько нужно изменить каждый вес. Если дисперсия градиентов \textbf{очень мала} (все градиенты близки к нулю), сеть почти не обучается — это называется \textit{затуханием градиентов} (vanishing gradients). Если дисперсия \textbf{слишком велика} (одни градиенты гигантские, другие крошечные), обучение становится хаотичным: веса могут «дёргаться» нескоординированно, что приводит к потере стабильности.

Идеальная ситуация — градиенты имеют умеренную, плавно уменьшающуюся дисперсию, что свидетельствует о равномерном обучении всех слоёв сети.

\subsubsection{Подробно}

Пусть модель имеет \(M\) параметров (весов). На каждой итерации вычисляется градиент функции потерь по каждому параметру: \(g_i = \frac{\partial \mathcal{L}}{\partial w_i}\). Вместо того чтобы анализировать все \(M\) скалярных градиентов, в данной работе используется упрощённый подход: вычисляются \textbf{нормы градиентов для каждого слоя} (или даже для каждого тензора параметров). Затем для списка этих норм \( \{ \|g^{(1)}\|, \|g^{(2)}\|, \dots, \|g^{(K)}\| \} \) вычисляется статистическая дисперсия:

\[
\text{GradVar} = \frac{1}{K}\sum_{k=1}^{K} \left( \|g^{(k)}\| - \overline{\|g\|} \right)^2,
\]
где \(\overline{\|g\|}\) — среднее арифметическое норм.

Такая оценка менее точна, чем дисперсия по всем отдельным параметрам, но значительно быстрее (не требует хранения миллионов значений) и даёт достаточное представление о разбросе «силы обновления» между разными частями сети.

\textbf{Интерпретация численных значений:}
\begin{itemize}
    \item \textbf{\(\text{GradVar} < 10^{-8}\)}: градиенты практически везде одинаково малы — скорее всего, сеть вошла в плато или затухание градиентов. Стоит проверить скорость обучения или архитектуру.
    \item \textbf{\(\text{GradVar} > 10^{-2}\)}: сильный разброс — некоторые слои обновляются гигантскими шагами, другие почти стоят на месте. Типично для начала обучения при слишком большой скорости обучения.
    \item \textbf{\(\text{GradVar}\) умеренная и постепенно убывает}: здоровое обучение, сеть адаптируется.
\end{itemize}

\subsubsection{Использование при обучении}

В данной работе GradVar вычисляется на \textbf{каждой итерации} (каждом шаге) внутри метода \texttt{TrainingLogger.log\_metrics}. Она сохраняется в CSV-файл (\texttt{train\_metrics.csv}) и может быть выведена в консоль или текстовый лог в соответствии с настройками YAML (флаги \texttt{include\_in\_file\_log} и \texttt{include\_in\_console}).

\textbf{Зачем нужна эта метрика?}
\begin{itemize}
    \item \textbf{Диагностика затухания градиентов:} особенно актуально для глубоких сетей, таких как NAFNet с 12 средними блоками. Если GradVar резко падает до нуля, это сигнал к изменению архитектуры (например, добавление skip-connections или нормализации).
    \item \textbf{Контроль стабильности:} слишком большая GradVar может указывать на необходимость снижения скорости обучения или градиентного клиппинга.
    \item \textbf{Мониторинг равномерности обучения:} если GradVar остаётся высокой на протяжении многих эпох, некоторые части сети могут недообучаться.
\end{itemize}

На практике GradVar используется исключительно для наблюдения и не влияет автоматически на процесс оптимизации. Однако её наличие в логах позволило в ходе экспериментов своевременно скорректировать гиперпараметры (например, снизить начальную скорость обучения с \(10^{-3}\) до \(10^{-4}\) при первых признаках нестабильности).

\subsection{Необходимая частота опроса}

\textbf{GradVar}, как и основные метрики (loss, PSNR), имеет смысл вычислять на \textbf{каждой итерации}. Причина: градиенты обновляются на каждом шаге, и их дисперсия может быстро меняться. Если вычислять GradVar только раз в несколько шагов, можно пропустить кратковременные всплески нестабильности, которые, тем не менее, способны нарушить обучение.

Однако на практике запись в CSV каждого значения GradVar (наряду с loss, PSNR и др.) не создаёт заметной нагрузки. В данном проекте частота опроса полностью совпадает с частотой логирования метрик, задаваемой параметром \texttt{file\_log\_freq} (по умолчанию 1, то есть каждый шаг). Вывод в консоль можно делать реже (например, каждые 10 шагов), чтобы не засорять экран, но для файлового лога рекомендуется сохранять каждое значение для последующего детального анализа.

\textbf{Итог:} GradVar должна вычисляться и сохраняться на каждом шаге (итерации) обучения.

\subsection{TV Ratio (Отношение полных вариаций)}

\subsubsection{Простое объяснение}

\textbf{TV Ratio} — это отношение «резкости» (или «гладкости») предсказанного изображения к эталонному (ground truth). Полная вариация (Total Variation, TV) — это сумма модулей разностей между соседними пикселями. Грубо говоря, TV измеряет, насколько сильно меняется яркость при переходе от пикселя к пикселю.

Представьте себе два изображения: \textbf{резкое} (с чёткими гранями) и \textbf{размытое} (с плавными переходами). У резкого изображения полная вариация будет \textit{большой}, потому что на границах соседние пиксели сильно отличаются. У размытого — \textit{маленькой}, так как значения соседних пикселей близки.

\textbf{TV Ratio} вычисляется как отношение TV предсказанного изображения к TV эталонного. Если TV Ratio близок к 1, значит, предсказание имеет такую же «шероховатость», как и эталон. Если TV Ratio < 1, предсказанное изображение \textit{гладче} эталона (потеряны высокие частоты, «мыло»). Если TV Ratio > 1, предсказание \textit{более шумное} или содержит лишние артефакты.

\subsubsection{Подробно}

Полная вариация для изображения \(I\) размера \(H \times W\) определяется как сумма модулей разностей по горизонтали и вертикали:
\[
\text{TV}(I) = \sum_{i=1}^{H-1}\sum_{j=1}^{W} |I(i+1,j) - I(i,j)| \;+\; \sum_{i=1}^{H}\sum_{j=1}^{W-1} |I(i,j+1) - I(i,j)|.
\]
В цветных изображениях (3 канала) TV обычно вычисляется для каждого канала отдельно и затем суммируется. Более эффективная для вычислений формула (используемая в коде):
\[
\text{TV}(I) = \sum_{i,j} |I(i+1,j) - I(i,j)| + \sum_{i,j} |I(i,j+1) - I(i,j)|.
\]

\textbf{TV Ratio} определяется как
\[
\text{TV Ratio} = \frac{\text{TV}(I_{\text{pred}})}{\text{TV}(I_{\text{gt}}) + \varepsilon},
\]
где \(\varepsilon\) — малая константа (например, \(10^{-8}\)), чтобы избежать деления на ноль. В коде проекта TV Ratio вычисляется на каждой итерации (если включено логирование) внутри \texttt{TrainingLogger.log\_metrics}.

\subsubsection{Использование при обучении}

\textbf{TV Ratio} — это не функция потерь, а вспомогательная метрика для контроля качества восстановления. Она помогает выявить две нежелательные тенденции:

\begin{itemize}
    \item \textbf{Чрезмерное сглаживание (TV Ratio < 0.9--1.0)}: модель выдаёт «мыльное» изображение — теряются мелкие детали, резкие грани размываются. Это характерно для обучения с агрессивным L2/L1-лоссом без частотной составляющей.
    \item \textbf{Избыточный шум и артефакты (TV Ratio > 1.1)}: модель генерирует лишнюю высокочастотную составляющую, похожую на шум или муар. Часто возникает при переобучении или слишком большом весе частотной потери.
\end{itemize}

В идеале TV Ratio должен быть близок к 1 и не сильно меняться в процессе обучения. В данной работе TV Ratio используется для тонкой настройки весов комбинированной потери: если TV Ratio заметно < 1, можно увеличить вес \(\gamma_2\) (FFL), чтобы модель активнее восстанавливала высокие частоты; если TV Ratio > 1, возможно, следует уменьшить \(\gamma_2\) или снизить скорость обучения.

\subsection{Необходимая частота опроса}

\textbf{TV Ratio} имеет смысл вычислять на \textbf{каждой итерации} (каждом шаге), аналогично loss и PSNR. Поскольку TV Ratio является чувствительным индикатором потери резкости или появления шума, его мониторинг с высоким временным разрешением позволяет быстро заметить нежелательные изменения в поведении модели.

В проекте TV Ratio вычисляется и логируется в CSV-файл с частотой, задаваемой параметром \texttt{file\_log\_freq} (обычно 1, т.е. каждый шаг). Вывод в консоль можно делать реже (через \texttt{print\_freq}), чтобы не перегружать экран. Рекомендуемая частота сохранения в файл — каждый шаг, поскольку объём данных невелик (одно число на итерацию), а для последующего анализа кривой TV Ratio полезно иметь все точки.

\textbf{Итог:} TV Ratio должен вычисляться и записываться на каждом шаге обучения (каждой итерации).

\subsection{SSIM (Structural Similarity Index)}

\subsubsection{Простое объяснение}

\textbf{SSIM} (индекс структурного сходства) — это метрика качества изображения, которая оценивает, насколько восстановленное изображение похоже на эталонное с \textbf{точки зрения человеческого восприятия}. В отличие от PSNR, который считает только пиксельные ошибки, SSIM учитывает три компонента:

\begin{itemize}
    \item \textbf{Яркость (luminance)} — среднюю интенсивность.
    \item \textbf{Контраст (contrast)} — разброс значений яркости.
    \item \textbf{Структуру (structure)} — пространственную организацию деталей (то, что для нас важнее всего).
\end{itemize}

Представьте, что вы сравниваете две фотографии одного и того же пейзажа. Если вы немного измените общую яркость (сделаете снимок светлее), PSNR может сильно упасть, а SSIM — остаться высоким, потому что структура (очертания гор, рек) сохранилась. Если же вы размоете границы (сотрёте структуру), SSIM сразу покажет низкое значение. Таким образом, SSIM гораздо лучше коррелирует с тем, что видит человек.

\subsubsection{Подробно}

Пусть \(x\) и \(y\) — два неперекрывающихся окна (патча) изображения (обычно \(11\times11\) пикселей). SSIM между ними определяется как

\[
\text{SSIM}(x,y) = \bigl[l(x,y)\bigr]^{\alpha} \cdot \bigl[c(x,y)\bigr]^{\beta} \cdot \bigl[s(x,y)\bigr]^{\gamma},
\]

где \(\alpha, \beta, \gamma\) — положительные веса (обычно \(\alpha=\beta=\gamma=1\)). Компоненты:

\begin{itemize}
    \item \textbf{Яркость:} \(l(x,y) = \frac{2\mu_x\mu_y + C_1}{\mu_x^2 + \mu_y^2 + C_1}\), где \(\mu_x\) — среднее яркости в окне.
    \item \textbf{Контраст:} \(c(x,y) = \frac{2\sigma_x\sigma_y + C_2}{\sigma_x^2 + \sigma_y^2 + C_2}\), где \(\sigma_x\) — стандартное отклонение.
    \item \textbf{Структура:} \(s(x,y) = \frac{\sigma_{xy} + C_3}{\sigma_x\sigma_y + C_3}\), где \(\sigma_{xy}\) — ковариация.
\end{itemize}

Константы \(C_1, C_2, C_3\) вводятся для устойчивости (чтобы не было деления на ноль). Итоговый SSIM для всего изображения получается как среднее арифметическое SSIM по всем окнам. Значение SSIM лежит в диапазоне \([-1, 1]\), но для изображений обычно близко к 1 (при полном сходстве). \(\text{SSIM} = 1\) означает идентичность.

\subsubsection{Использование при обучении}

В данной работе \textbf{SSIM используется только как оценочная метрика} (не участвует в оптимизации). Она вычисляется во время валидации на тестовых изображениях. В отличие от PSNR, SSIM показывает, насколько хорошо модель восстанавливает структуру, а не просто пиксельную точность. Для задачи демозаики это особенно важно, потому что муар и артефакты могут существенно искажать структуру даже при приемлемом PSNR.

В диссертации SSIM используется для:
\begin{itemize}
    \item \textbf{Сравнения с классическими методами} (билинейная интерполяция, MHC).
    \item \textbf{Выбора лучшей модели} (по валидационному SSIM).
    \item \textbf{Визуального контроля} (сохранённые в валидации изображения сопровождаются значениями SSIM).
\end{itemize}

\textbf{Примечание:} в текущей реализации SSIM вычисляется только на валидации, а не на каждом шаге. Это связано с тем, что расчёт SSIM относительно затратен, а его обновление раз в несколько эпох даёт достаточное представление о качестве.

\subsection{Необходимая частота опроса}

\textbf{SSIM} имеет смысл вычислять \textbf{на валидации} — то есть каждые \(k\) эпох (в данном проекте параметр \texttt{validation\_freq} задаётся в конфиге, обычно 1 или 5). Нет необходимости вычислять SSIM на каждом шаге (это было бы слишком медленно и не дало бы полезной информации, так как изменения от батча к батчу незначительны).

\textbf{Рекомендуемая частота:} каждую эпоху (или каждые 5 эпох для более быстрых экспериментов). Значения SSIM записываются в лог (через \texttt{logger.info}) и могут быть извлечены для построения графика с помощью скрипта \texttt{plot\_metrics.py}.

\textbf{Итог:} SSIM вычисляется редко (по эпохам) и служит для общей оценки качества, а не для контроля сходимости в реальном времени. Это соответствует типичной практике в задачах реставрации изображений.

\subsection{Использование видеопамяти (VRAM Usage)}

\subsubsection{Простое объяснение}

\textbf{VRAM} (Video RAM) — это память, расположенная непосредственно на видеокарте (GPU). В отличие от обычной оперативной памяти (RAM), VRAM имеет очень высокую пропускную способность и предназначена для хранения данных, с которыми GPU работает непосредственно: веса нейросети, промежуточные активации, градиенты, входные тензоры и т.д.

В процессе обучения глубоких моделей (таких как NAFNet) контроль за использованием VRAM критически важен, потому что при превышении доступного объёма возникает ошибка \texttt{Out of Memory (OOM)}. Тогда обучение прерывается, а все потраченные часы могут оказаться напрасными. Мониторинг VRAM позволяет вовремя уменьшить размер батча, оптимизировать архитектуру или использовать смешанную точность (FP16).

\subsubsection{Подробно}

В проекте фиксируются две основные величины:

\begin{itemize}
    \item \textbf{Allocated VRAM} — объём памяти, фактически выделенный PyTorch под тензоры и графы вычислений (то, что реально используется).
    \item \textbf{Reserved VRAM} — память, зарезервированная кэширующим аллокатором PyTorch (может быть больше выделенной, чтобы ускорить дальнейшие запросы).
\end{itemize}

Обе величины измеряются в гигабайтах (ГБ). В коде (метод \texttt{log\_metrics} в \texttt{TrainingLogger}) они получаются через
\begin{verbatim}
torch.cuda.memory_allocated(device) / (1024**3)
torch.cuda.memory_reserved(device) / (1024**3)
\end{verbatim}

Разница между \texttt{Reserved} и \texttt{Allocated} обычно невелика и указывает на эффективность работы аллокатора. Если она становится слишком большой (например, зарезервировано 10 ГБ, а используется 2 ГБ), это может свидетельствовать о фрагментации памяти или неоптимальных настройках.

\subsubsection{Использование при обучении}

\textbf{VRAM Usage} не влияет на процесс оптимизации, но является важной \textit{диагностической метрикой}. Она помогает:

\begin{itemize}
    \item \textbf{Избежать OOM-ошибок:} если в логах видно, что зарезервированная память приближается к физическому объёму GPU (например, 11.5 ГБ из 12), необходимо срочно уменьшить \texttt{batch\_size} или разрешение входных патчей.
    \item \textbf{Оптимизировать использование ресурсов:} при обучении на нескольких GPU метрика позволяет балансировать нагрузку.
    \item \textbf{Обнаружить утечки памяти:} если выделенная память неуклонно растёт с каждой итерацией, это признак того, что где-то сохраняются тензоры, которые не удаляются.
\end{itemize}

В данной работе VRAM Usage логируется в CSV-файл на каждой итерации (каждом шаге) и выводится в консоль (если включено в \texttt{include\_in\_console}). Это позволяет в реальном времени следить за расходом памяти.

\subsection{Необходимая частота опроса}

\textbf{VRAM Usage} имеет смысл измерять \textbf{на каждой итерации} (каждом шаге). Память может меняться непредсказуемо: например, при создании временных тензоров внутри модели или при смене этапа (тренировка/валидация). Разовые замеры могут пропустить пики, которые приводят к OOM.

В проекте частота опроса полностью совпадает с основной частотой логирования (параметр \texttt{file\_log\_freq}, по умолчанию 1). Вывод в консоль регулируется отдельно (обычно каждые 10--20 шагов, чтобы не засорять экран). Для файлового CSV-лога рекомендуется сохранять каждое значение, так как объём данных невелик, а для последующего анализа пиков VRAM полезно иметь все точки.

\textbf{Итог:} VRAM Usage должен вычисляться и записываться на каждом шаге обучения (итерации).

\subsection{Гиперпараметры обучения и их влияние}

В данном разделе описываются параметры, которые непосредственно управляют процессом обучения нейросети NAFNet: размер батча, количество эпох, тип оптимизатора и планировщика, коэффициенты регуляризации и веса компонентов комбинированной потери. Эти значения не выводятся из данных, а задаются исследователем (или определяются из опыта) и оказывают решающее влияние на скорость сходимости, устойчивость и конечное качество модели.

\subsubsection{Размер батча (Batch size)}

\textbf{Простое объяснение:}  
Батч — это количество изображений (патчей), которые обрабатываются сетью одновременно перед обновлением весов. Представьте, что вы учитесь решать примеры: если просматривать по одному примеру и сразу корректировать свои знания, это быстро, но шумно (каждый пример может сбивать с толку). Если просматривать сразу много примеров и усреднять их, корректировка будет более надёжной, но потребует больше памяти. Размер батча — это компромисс между стабильностью обновлений и требованиями к видеопамяти.

\textbf{Подробно:}  
Пусть \(B\) — размер батча. Градиент функции потерь вычисляется как среднее по \(B\) примерам: \(\nabla \mathcal{L} = \frac{1}{B}\sum_{i=1}^{B} \nabla \mathcal{L}_i\). Чем больше \(B\), тем точнее оценка градиента (меньше дисперсия), но тем больше памяти требуется для хранения промежуточных активаций. Слишком маленький батч (например, \(B=1\)) приводит к хаотичным обновлениям и может ухудшить сходимость. Слишком большой батч может вызвать \texttt{Out of Memory} и не даёт значимого выигрыша в качестве после определённого порога. В данной работе выбран \(B = 16\) (настраивается в YAML как \texttt{batch\_size\_per\_gpu}). Это значение позволяет полностью загрузить память GPU (12 ГБ RTX 4070 Ti) при размере патча \(256\times256\) и 4 входных каналах, обеспечивая стабильные градиенты.

\textbf{Использование при обучении:}  
Размер батча задаётся в конфигурационном файле и не меняется в процессе обучения. Он влияет на скорость обучения (чем больше батч, тем можно использовать больший LR). В данном проекте LR = \(10^{-3}\) подобран под \(B=16\).

\subsubsection{Количество эпох (Number of epochs)}

\textbf{Простое объяснение:}  
Эпоха — это один полный проход всего обучающего набора данных через сеть. Если датасет содержит 1000 патчей, а батч — 16, то одна эпоха состоит из \(1000/16 \approx 63\) шагов. Количество эпох определяет, сколько раз сеть «увидит» каждый пример. Слишком мало эпох — сеть недообучится, слишком много — может переобучиться (запомнить шум).

\textbf{Подробно:}  
В данной работе установлено \texttt{num\_epochs = 150} (конфигурационный параметр). Это значение было выбрано эмпирически: к 100–150 эпохам значения loss и PSNR стабилизируются на плато, а дальнейшее обучение приводит лишь к незначительному улучшению (менее 0.5 дБ). При увеличении числа эпох до 300 переобучения не наблюдалось из-за малого размера датасета и применения аугментаций.

\textbf{Использование при обучении:}  
Количество эпох задаётся до старта и определяет длительность обучения. Параметр \texttt{save\_checkpoint\_epoch} позволяет сохранять промежуточные веса каждые \(k\) эпох, что удобно для анализа.

\subsubsection{Оптимизатор (AdamW)}

\textbf{Простое объяснение:}  
AdamW — это алгоритм, который обновляет веса сети. Он комбинирует идеи адаптивного градиента (разная скорость для каждого веса) и импульса (учёт предыдущих обновлений). Это один из самых надёжных оптимизаторов для задач компьютерного зрения. «W» означает, что регуляризация весов (weight decay) применяется отдельно, что часто даёт лучшие результаты, чем в стандартном Adam.

\textbf{Подробно:}  
AdamW использует экспоненциально скользящие средние градиента и квадрата градиента:
\[
m_t = \beta_1 m_{t-1} + (1-\beta_1) g_t, \quad v_t = \beta_2 v_{t-1} + (1-\beta_2) g_t^2,
\]
затем корректирует смещения и обновляет веса:
\[
w_{t+1} = w_t - \eta \left( \frac{\hat{m}_t}{\sqrt{\hat{v}_t}+\varepsilon} + \lambda w_t \right),
\]
где \(\lambda\) — коэффициент weight decay. В данной работе \(\beta_1 = 0.9\), \(\beta_2 = 0.9\), \(\varepsilon = 10^{-8}\), \(\lambda = 0\). Выбор \(\lambda=0\) объясняется тем, что модель не склонна к переобучению на малом датасете, а также для чистоты эксперимента (регуляризация может подавлять высокие частоты).

\textbf{Использование при обучении:}  
AdamW инициализируется в начале и работает на всех итерациях. Его параметры не изменяются в процессе, кроме глобальной скорости обучения, которой управляет планировщик.

\subsubsection{Планировщик скорости обучения (CosineAnnealingLR)}

\textbf{Простое объяснение:}  
Этот планировщик плавно снижает скорость обучения от начального значения до минимального по косинусоиде. Вместо резких ступенек (как у StepLR) он обеспечивает мягкое уменьшение, что помогает сети аккуратно «дойти до дна» функции потерь.

\textbf{Подробно:}  
Формула: \(\eta(t) = \eta_{\min} + \frac{1}{2}(\eta_{\max} - \eta_{\min})(1 + \cos(\frac{t}{T_{\max}}\pi))\), где \(t\) — номер эпохи (или шага, но в проекте используется по эпохам). Начальное значение \(\eta_{\max} = 10^{-3}\), конечное \(\eta_{\min} = 10^{-7}\), \(T_{\max} = 150\) (число эпох). Такое сильное уменьшение гарантирует, что в конце обучения сеть делает очень маленькие шаги, тонко подстраивая веса.

\textbf{Использование при обучении:}  
Планировщик вызывается после каждой эпохи (\texttt{scheduler.step()}). Его влияние логируется в CSV.

\subsubsection{Веса компонентов комбинированной потери (\(\gamma_1, \gamma_2\))}

\textbf{Простое объяснение:}  
Комбинированная потеря \(\mathcal{L}_{\text{total}} = \gamma_1 \cdot \mathcal{L}_1 + \gamma_2 \cdot \mathcal{L}_{\text{FFL}}\) объединяет пространственный (L1) и частотный (FFL) критерии. Веса \(\gamma_1, \gamma_2\) определяют, какой аспект важнее. Если \(\gamma_2\) слишком мал, сеть будет игнорировать высокие частоты и давать «мыло». Если \(\gamma_2\) слишком велик, возможны шумы и артефакты.

\textbf{Подробно:}  
В данном проекте установлены \(\gamma_1 = 1.0\), \(\gamma_2 = 1.0\) (равный вклад). Это значение было получено после нескольких пробных запусков: при \(\gamma_2 = 0.5\) резкость была недостаточной, при \(\gamma_2 = 2.0\) появлялись цветные артефакты. Параметры задаются в YAML-конфиге в секции \texttt{losses}.

\textbf{Использование при обучении:}  
Веса не изменяются в процессе. Они влияют на каждую итерацию при вычислении loss и градиентов.

\subsubsection{Другие гиперпараметры (неиспользуемые)}

В данном проекте \textbf{не используются} градиентное клиппинг (gradient clipping) и L2-регуляризация (\texttt{weight\_decay}=0). Это сознательный выбор, упрощающий эксперимент и не влияющий на основные выводы.

\subsection{Резюме по гиперпараметрам}

Все перечисленные параметры (размер батча, количество эпох, оптимизатор, планировщик, веса лосса) задаются в конфигурационном файле \texttt{RAW\_NAFNet\_NikonD600.yml} в соответствующих секциях. Они не выводятся из данных, а определяются исходя из архитектуры сети, объёма видеопамяти и экспериментального подбора. Выбранные значения обеспечивают стабильную сходимость за приемлемое время (около 6--8 часов на RTX 4070 Ti) и дают воспроизводимые результаты.

\end{document}


\end{document}